\section{Hilbert Spaces} 

	\begin{definition}[Hilbert Space]
	  \label{def:hilbert-space}
		A \textbf{Hilbert space} is a complete \hyperref[lin-def:inner-product]{inner product} space. 
	\end{definition}

	In analysis, you want to solve problems by construting sequences. It's bad if your \hyperref[lin-def:vector-space]{vector space} is not complete, i.e. these sequences don't converge. 

	\begin{lemma}[Nearest Point in a Closed Convex Subset of a Hilbert Space]
	  \label{thm:nearest-point-in-a-closed-convex-subset-of-a}
	  Let $K$ be a closed, nonempty convex subset of \hyperref[def:hilbert-space]{Hilbert space} $H$. Then, there exists a unique element of minimal norm in $K$. 
	\end{lemma}
	\begin{proof}
		You should remember this proof since it's used in the \hyperref[thm:riesz-representation-theorem]{Riesz Representation theorem}'s proof. Let $\delta = \inf_{x \in K} |x|$. Pick a sequence $(x_j) \subset K$ s.t. $|x_j| \to \delta$ as $j \to +\infty$. We claim that $(x_j)$ is Cauchy. This is somewhat nonobvious, so as a proof, let $\epsilon > 0$. Then, there exists a $N \in \mathbb{N}$ s.t. $j > N \implies |x_j|^2 \leq \delta^2 + \frac{\epsilon^2}{4}$. So by convexity, 
		\begin{equation}
			m, n > N \implies \frac{x_m + x_n}{2} \in K
		\end{equation}
		Since $\delta$ is an infimum, 
		\begin{equation}
			\bigg| \frac{x_m + x_n}{2} \bigg| \geq \delta \implies |x_m + x_n|^2 \geq 4 \delta^2
		\end{equation}
		By expanding the quadratic 
		\begin{align}
			|x_m - x_n|^2 & = |x_m|^2 + |x_n|^2 - 2 \langle x_m, x_n \rangle \\ 
										& = - |x_m + x_n|^2 + 2|x_m|^2 + 2|x_n|^2 \\ 
										& < - 4 \delta^2 + 4 \delta^2 + 4 \frac{\epsilon^2}{4} = \epsilon^2
		\end{align} 
		which implies that $|x_m - x_n| < \epsilon$ if $m, n > N$. So, it is Cauchy\footnote{Check if this is related to strongly Cauchy?} and since $H$ is complete, $(x_j)$ has a limit---call it $x_\infty \in K$ (since $K$ is closed) with $|x_\infty| = \delta$. Now we prove uniqueness. Suppose $y \in K$ satisfies $|y| = \delta \implies \frac{x_\infty + y}{2} \in K \implies |x_\infty + y|^2 > 4 \delta^2$. So, $|x_\infty - y|^2 = 0 \implies x_\infty = y$. 
	\end{proof}
    
	\begin{theorem}[Riesz Representation Theorem]
	  \label{thm:lp-duality}
	  Let $H$ be a \hyperref[def:hilbert-space]{Hilbert space}. Define $\Lambda: H \to H^\ast$ by $\Lambda(x)(y) \coloneqq \langle y, x \rangle$. Then, $\Lambda$ is an isometeric isomorphism. 
	\end{theorem}
	\begin{proof}
	  For $x \neq 0$, we have 
		\begin{equation}
			\|\Lambda (x)\| = \sup_{y \neq 0} \frac{|\Lambda(x)(y)|}{|y|} = \sup_{y \neq 0} \frac{|\langle y, x \rangle|}{|y|} \leq \frac{|x||y|} {|y|} = |x|
		\end{equation}
		For the other direction, 
		\begin{equation}
			\frac{|\Lambda(x)(x)|}{|x|} = \frac{|x|^2}{|x|} = |x|
		\end{equation}
		Therefore, $\|\Lambda(x)(x)\| = |x|$, and $\Lambda$ is an isometry. 
		\begin{enumerate}
		  \item \textit{Injective}. It is easy to show that $\Lambda: H \to H^\ast$ is injective. 
			\item \textit{Surjective}. This is more involved, and we need convexity. For simplicity, let our field be $\mathbb{R}$. We've shown that $\Lambda: H \to H^\ast$ is an isometric injection. Let $\phi \in H^\ast \setminus \{0\}$. There are 2 families of natural \hyperref[st-def:convex-sets]{convex sets}: affine spaces and balls. It is easy to see that $K = \phi^{-1} (1)$ is an affine space and hence convex. Let $x_0$ be an element of minimal norm of $K$. We claim that $x_0 \perp \mathrm{ker}(\phi)$. Indeed, we have 
				\begin{align}
					y \in \mathrm{ker}(\phi) & \implies x_0 + ty \in K \text{ for all } t \in \mathbb{R} \\ 
																	 & \implies |x_0|^2 \leq |x_0 + ty|^2 = |x_0|^2 + 2t \langle x_0, y \rangle + t^2 |y|^2 \\ 
																	 & \implies 0 \leq 2t \langle x_0, y \rangle + t^2 |y|^2
				\end{align}
				Take $t = -\frac{\langle x_0, y \rangle}{|y|^2}$, then we have 
				\begin{equation}
					0 \leq \frac{-2 \langle x_0, y \rangle^2}{|y|^2} + \frac{\langle x_0, y \rangle^2}{|y|^2} \implies 0 \leq - \frac{\langle x_0, y \rangle^2}{|y|^2} 
				\end{equation}
				and so $\langle x_0, y \rangle = 0$. So to summarize, $\phi(x_0) = 1$ with $x_0 \perp \mathrm{ker}(\phi)$. Let $x \in H$. Then, 
				\begin{equation}
					\phi\big( x - \phi(x) x_0 \big) = \phi(x) - \phi(x) \underbrace{\phi(x_0)}_{=1} = 0
				\end{equation}
				But this means 
				\begin{align}
					x - \phi (x) x_0 \in \mathrm{ker}(\phi) 
						& \implies 0 = \langle x - \phi(x) x_0, x_0 \rangle \\ 
						& \implies \langle x, x_0 \rangle = \phi(x) \langle x_0, x_0 \rangle = \phi(x) |x_0|^2 \\ 
						& \implies \phi(x) = \bigg\langle x, \frac{x_0}{|x_0|^2} \bigg\rangle \\ 
						& \implies \phi = \Lambda \bigg( \frac{x_0}{|x_0|^2} \bigg)
				\end{align}
				and therefore, $\Lambda$ is surjective. 
		\end{enumerate} 
	\end{proof}

	\begin{definition}[Orthogonal Subspace]
	  \label{def:orthogonal-subspace}
	  Let $S \subset H$ be a subspace of a \hyperref[def:hilbert-space]{Hilbert space}, and let 
		\begin{equation}
			S^\perp \coloneqq \{x \in H \colon \langle s, x \rangle = 0 \forall s \in S\}
		\end{equation} 
	\end{definition}


	\begin{theorem}[1.4.6]
	  \label{thm:1-4-6}
	  Let $E \subset H$ be a closed subspace of \hyperref[def:hilbert-space]{Hilbert space} $H$. Then, $H = E \oplus E^\perp$. 
	\end{theorem}
	\begin{proof}
	  If $x \in E \cap E^\perp$, then $x \perp x \iff \langle x, x = 0 \iff x = 0$. If $x \in H$, then $x + E$ is a closed \hyperref[mv-def:convex-set]{convex set}, which implies that there exists an element $x_0$ of least norm in $x + E$. Therefore, $x_0 \perp E$ as before (previous argument: $|x_0 + ty|^2 \geq |x_0|^2$...), which implies that 
		\begin{equation}
			x = \underbrace{x - x_0}_{\in E} + \underbrace{x_0}_{\in E^\perp}
		\end{equation}
	\end{proof}
